% Half-title, title and copyright pages.
%
% \bookversion is the release the PDF was built for: `make VERSION=book-v2` passes it in, and the
% release workflow passes the tag. Without it the title page carries only the date of the build.
\providecommand{\bookversion}{}
\newcommand{\bookdate}{\number\day\space\ifcase\month\or January\or February\or March\or April\or
  May\or June\or July\or August\or September\or October\or November\or December\fi\space
  \number\year}
\newcommand{\bookedition}{\ifx\bookversion\empty\else\bookversion\,·\,\fi\bookdate}

\thispagestyle{empty}
\vspace*{2.2in}
\begin{flushright}
  {\Large\bfseries Programmable Logic}
\end{flushright}
\cleardoublepage

\thispagestyle{empty}
\vspace*{1.1in}
\begin{flushleft}
  {\color{accent}\rule{0.9in}{2.5pt}}\par\vspace{1.4ex}
  {\fontsize{32}{38}\selectfont\bfseries Programmable\\[0.15ex] Logic\par}
  \vspace{2.2ex}
  {\Large\itshape Digital design in VHDL, from a gate\\network to a CAN node in an FPGA\par}
  \vspace{1.2in}
  {\large Erik Pihl\par}
  \vfill
  {\small\scshape\lsstyle qrtech academy\par}
  \vspace{0.6ex}
  {\small\color{muted} \bookedition\par}
\end{flushleft}
\newpage

\thispagestyle{empty}
\vspace*{\fill}
{\small
\noindent\textit{Programmable Logic}\\
Copyright \textcopyright\ 2026 QRTECH Academy.

\medskip
\noindent This is the English edition of \textit{Programmerbar logik}. The book is typeset from the
course's own material: twenty lectures with their appendices, their exercises, the worked examples,
the group project's specification and the two documents that form the contract with the driver
course. The code in the book, and all the code in the accompanying course repository, may be used
under the MIT licence, which is the repository's licence:

\medskip
\noindent\url{https://opensource.org/licenses/MIT}

\medskip
\noindent The course repository holds the lecture material, the worked examples, every testbench
handed out, the build scripts and the Python sources behind the figures. Every path the book refers
to, for example \file{lectures/L06/serial_rx8/serial_rx8.vhd}, is a path in that repository.

\medskip
\noindent The VHDL code in the book is VHDL-93, simulated with GHDL and synthesised with Quartus
Prime Lite. The C++ code in \secref{c10:sec:framing} is C++17, compiled with GCC. DE0-CV,
Cyclone~V, Quartus, USB-Blaster, AVR and AVR32DB28 are trademarks belonging to their respective
owners. CAN is standardised in ISO~11898; the controller this book builds is a deliberately
simplified teaching design and no certifiable implementation, and \secref{c19:sec:gaps} lists the
differences.

\medskip
\noindent Text and figures are in English. So are all the code, all the comments in the code and
all the output from tools and testbenches, exactly as in the course material.

\medskip
\noindent Typeset in TeX Gyre Pagella and DejaVu Sans Mono with Lua\LaTeX.

\medskip
\noindent\textcolor{muted}{This edition: \bookedition.}
\par}
\cleardoublepage
